\documentclass[11pt]{article}
\usepackage[T1]{fontenc}\usepackage[utf8]{inputenc}
\usepackage{lmodern,amsmath,amssymb,amsthm,mathtools,microtype,xcolor}
\usepackage[margin=1in]{geometry}
\usepackage[colorlinks=true,linkcolor=blue!50!black,urlcolor=blue!50!black,citecolor=blue!50!black]{hyperref}
\newtheorem{theorem}{Theorem}[section]\newtheorem{lemma}[theorem]{Lemma}
\newtheorem{proposition}[theorem]{Proposition}\newtheorem{corollary}[theorem]{Corollary}
\theoremstyle{remark}\newtheorem{remark}[theorem]{Remark}
\newcommand{\R}{\mathbb R}\newcommand{\E}{\mathbb E}\newcommand{\Prb}{\mathbb P}
\newcommand{\Lip}{\mathrm{Lip}}\newcommand{\sech}{\mathrm{sech}}
\newcommand{\norm}[1]{\left\Vert#1\right\Vert}
\setlength{\emergencystretch}{3em}
\hypersetup{pdftitle={Exponential depth dependence in Bayesian neural-network regression},pdfauthor={ChatGPT-6 Astra; Kaarel Hanni; Dmitry Vaintrob}}
\title{Exponential depth dependence in Bayesian neural-network regression}
\author{ChatGPT-6 Astra \and Kaarel H\"anni \and Dmitry Vaintrob\thanks{Authors are listed in alphabetical order. The main learning statements have Lean proofs; see the repository for scope and exclusions.}}
\date{October 18, 2026}
\begin{document}\maketitle
\begin{abstract}
We construct bounded-weight tanh teachers of constant width and depth $L$
for which a Gaussian Bayesian neural-network learner requires
$\exp(\Omega(L))$ examples to generalize. Inputs follow one fixed nonatomic
distribution on an arc of the unit circle, and training observes the
teacher's real-valued outputs. The posterior conditions only on fitting
these outputs within a fixed positive tolerance. In the strong-prior
limit, and at every finite prior exponent $\gamma>1$ at widths above a
polynomial threshold in sample size and depth, the posterior fails to assign probability greater than
$2/3$ to a correct prediction on at least half the test inputs, with sample
probability at least $49/50$.
The teacher iterates a fixed expanding map. Its outputs differ by a constant
on exponentially close paired inputs. A sample of subexponential size,
however, admits an exact interpolant with sufficiently low quadratic
energy that the posterior concentrates on functions, or uniform
approximations to functions, having much smaller Lipschitz constants.
The pairing argument handles arbitrary posterior mixtures. All
minimum-energy exact interpolants fail as well, including the specified
iterated strong-prior, zero-tolerance posterior limits.
Thus, for these learners, no sample bound that holds uniformly over
nonatomic input distributions can be jointly polynomial in teacher width,
weight bound, and depth. The fixed distribution is a Cantor measure; allowing the
distribution to vary with depth also gives smooth arclength densities.
Neither construction uses uniform measure on the circle.
A separate width-first argument gives the same exponential obstruction
at the mean-field endpoint $\gamma=1$, for the same fixed input law,
without a quantified width threshold.
The original Boolean-cube classification problem remains outside
these results.
\end{abstract}

\section{Introduction}
Can a Bayesian neural network learn every small bounded-weight teacher
with a number of examples polynomial in both width and depth, without
restricting the networks in its prior? On arbitrary continuous input
distributions, the answer is negative for the Gaussian learners considered
here. We give a family with fixed input dimension, fixed teacher width,
fixed weight bound, fixed training and testing tolerances, and a single
fixed input distribution. Only the teacher depth grows. Learning requires
exponentially many examples in that depth.

The obstruction is a property of the learning rule, rather than an
information-theoretic impossibility for every algorithm. A teacher can
iterate a small expanding map, producing a highly oscillatory function.
The independent Gaussian weight prior instead rewards a low-energy
interpolant of the observed sample. At large learner width, such an
interpolant can transmit a scalar signal through many layers with a
small excess energy cost. Its product of layer norms, and hence its
Lipschitz constant, is polynomial in the sample size. Below an exponential
sample size this smoothness prevents it from approximating the teacher.

The proof requires more than a large Lipschitz constant for one teacher
implementation. We prove an approximation barrier for \emph{every}
function of small Lipschitz constant, exhibit a low-energy competitor
for each typical sample, and transfer the resulting competition to the
full posterior. A measure-preserving pairing of test inputs shows that
mixing individually inaccurate predictions cannot repair the confidence
guarantee.

We prove the strong-prior result first, then give a finite-scale extension
for every fixed exponent above one. The latter uses Gaussian isoperimetry
to remove a typical small-operator-norm Gaussian component. The actual
posterior parameter vector need not itself have small Frobenius energy.
We include the required memorization and perturbation arguments, so the
proof does not require another paper from this research program.
At the mean-field endpoint, a separate argument uses entropy costs and
finite-dimensional large deviations. It establishes the same depth
obstruction for the fixed nonatomic distribution as width tends to
infinity, without a quantitative threshold on width.

\section{Model and main conclusions}\label{sec:main}
For input $x\in S^1\subset\R^2$, write $\bar x=(1,x)$ and consider
depth-$L$ networks with $L-1$ hidden layers, all of width $N$:
\[
 h^0(x)=\bar x,\qquad
 h^j(x)=\tanh(W^jh^{j-1}(x))\quad(1\le j<L),\qquad
 z_w(x)=v^Th^{L-1}(x).
\]
There are no separate internal biases. More generally, in input dimension
$d$ the independent centered Gaussian prior has variances
\begin{equation}\label{eq:scales}
 \mathrm{Var}(W^1_{ij})=\sigma^2N/d,\qquad
 \mathrm{Var}(W^j_{ab})=\sigma^2\ (2\le j<L),\qquad
 \mathrm{Var}(v_a)=\sigma^2/N.
\end{equation}
Its density in energy coordinates is proportional to
$\exp[-\kappa_N(w)/(2\sigma^2)]$, where
\begin{equation}\label{eq:energy}
 \kappa_N(w)=\frac dN\norm{W^1}_F^2+
 \sum_{j=2}^{L-1}\norm{W^j}_F^2+N\norm v_2^2.
\end{equation}
The strong-prior limit means $\sigma\downarrow0$ at fixed finite $N,L$.
The finite exponent family sets $\sigma^2=N^{-\gamma}$.

For a target $f$, independent sample inputs $X_i\sim D$, and observed
real labels $Y_i=f(X_i)$, let $\Pi_{S,\sigma}$ be this prior conditioned on
\begin{equation}\label{eq:training}
 |z_w(X_i)-Y_i|<\rho\qquad(1\le i\le n).
\end{equation}
There is no condition on unobserved inputs and no global output margin.
Throughout the main results we fix
\[
 \rho=1/8,\qquad \tau=1/4,
\]
and define the posterior's correct-output probability by
\[
 q_{S,\sigma}(x)=\Pi_{S,\sigma}\{|z_w(x)-f(x)|\le\tau\}.
\]
Thus an input is a confidence error if $q_{S,\sigma}(x)\le2/3$.

Here are explicit sample-dependent resource bounds used in the theorems:
\begin{equation}\label{eq:resources}
 \alpha=\frac{\log2}{\log20},\quad
 \Delta_n=(100n^2)^{-1/\alpha},\quad
 a_n=\Delta_n^{-1}\log(16n),\quad R_n=\max\{2,a_n\},
\end{equation}
\begin{equation}\label{eq:MT}
 M_n=\left\lceil\max\{4,n,6nR_n^2\}\right\rceil,
 \qquad T_n=4M_n n,\qquad
 K_n(L)=2+(L+14)T_n^{1/(L-1)}.
\end{equation}
These bounds are fixed polynomials in $n$ up to logarithmic factors.
In particular $M_n$ and $T_n$ admit universal polynomial upper bounds.

\begin{theorem}[Exponential lower bound in the strong-prior limit]
\label{thm:strong}
There exist a fixed nonatomic Borel probability distribution $D$ on
$S^1$, targets $f_r:S^1\to\R$ for all integers $r\ge1$, and a universal
integer $r_0$ with the following properties. Each $f_r$ is implemented
by a tanh teacher of width three, depth $L=r+1$, and weights of absolute
value at most $30$. It satisfies $|f_r|\le13/20$ on the support of $D$.
For every $r\ge r_0$, every integer $1\le n\le e^{r/100}$, learner depth
$L=r+1$, and every fixed finite learner width $N\ge\max\{30,M_n\}$,
\begin{equation}\label{eq:strong-result}
 \liminf_{\sigma\downarrow0}\Prb_{S\sim D^n}
 \left\{D\{x:q_{S,\sigma}(x)\le2/3\}\ge1/2\right\}
 \ge49/50.
\end{equation}
The distribution, input dimension, teacher width and weight bound, and
both tolerances are independent of $r,n,N$.
\end{theorem}

Padding the teacher to width $30$ makes its weights bounded by its width,
if that is the chosen definition of teacher complexity. Thus its
width-based complexity remains an absolute constant. The learner widths
in Theorem~\ref{thm:strong} are allowed by any guarantee asserted for
\emph{every} width above teacher complexity. The theorem does not claim
failure at every smaller width, including width equal to the teacher's.

\begin{theorem}[Exponential lower bound at finite exponents above one]
\label{thm:finite}
Use the distribution, targets and constants of
Theorem~\ref{thm:strong}. Fix $\zeta>0$. There is a width threshold
$\mathcal N_\zeta(n,L)$ with a polynomial upper bound in $n,L$, whose
degree may depend on $\zeta$, such that for every $r\ge r_0$,
$1\le n\le e^{r/100}$, $L=r+1$, every
$N\ge\mathcal N_\zeta(n,L)$, and every
$0<\sigma^2\le N^{-(1+\zeta)}$,
\[
 \Prb_{S\sim D^n}
 \left\{D\{x:q_{S,\sigma}(x)\le2/3\}\ge1/2\right\}\ge49/50.
\]
In particular this applies at the actual finite scale
$\sigma^2=N^{-\gamma}$ for each fixed $\gamma>1$, by taking
$\zeta=\gamma-1$. An explicit threshold is given in
Lemma~\ref{lem:finite-ceiling}.
\end{theorem}

The same construction obstructs minimum-energy exact interpolation and
a specified iterated Bayesian limit; see Section~\ref{sec:exact}.
At $\gamma=1$, Theorem~\ref{thm:critical-regression-depth-lower}
gives a separate width-first lower bound: on a sample event of probability
at least $49/50$, the lower limit of the confidence-error mass as
$N\to\infty$ is at least $1/2$, for an exponential range of sample sizes.
This endpoint theorem does not quantify how large the width must be.
These conclusions are specific to the stated learner. They are not a
lower bound for every method that knows the teacher family.

\section{A constant-width expanding teacher}\label{sec:map}
Define
\begin{equation}\label{eq:map}
 J=[-3/4,3/4],\qquad
 F(t)=2\{\tanh(10(t+1/2))-\tanh(10(t-1/2))\}-1.
\end{equation}

\begin{lemma}[Two expanding branches]\label{lem:branches}
There are inverse branches $\phi_\pm:J\to J$ of $F$, with images
in $[1/2,13/20]$ and $[-13/20,-1/2]$, respectively, such that
\begin{equation}\label{eq:branches}
 \frac1{20}|s-t|\le|\phi_\pm(s)-\phi_\pm(t)|
 \le\frac13|s-t|\qquad(s,t\in J).
\end{equation}
The two images have distance at least one.
\end{lemma}
\begin{proof}
The map $F$ is even. On $[1/2,13/20]$ it is strictly decreasing, and
\[
 F(1/2)=2\tanh(10)-1>3/4,
 \qquad F(13/20)<1-2\tanh(3/2)<-3/4.
\]
Moreover
\[
 20\{\sech^2(3/2)-\sech^2(10)\}
 \le |F'(t)|\le20.
\]
For explicit elementary bounds, the degree-four Taylor sum gives
$e^3>15$, while the Taylor series and its geometric remainder give
$e<11/4$, hence $e^3<21$. These inequalities imply
$\tanh(3/2)>7/8$ and $\sech^2(3/2)>1/6$, while
$e^{20}>20^4/24>4000$ gives $\sech^2(10)<1/1000$.
Thus the lower derivative bound exceeds three. The inverse function
theorem gives the positive branch and its derivative bounds. Reflection
gives the negative branch. Their image locations give the gap.
\end{proof}

\begin{lemma}[Teacher implementation]\label{lem:teacher}
For every $r\ge1$, the function $f_r(x)=F^r(x_1)$ on $S^1$ has a
width-three, depth-$(r+1)$ tanh implementation with weights bounded by $30$.
\end{lemma}
\begin{proof}
The first hidden layer computes
\[
 a_1=\tanh(10x_1+5),\qquad
 a_2=\tanh(10x_1-5),\qquad c=\tanh(\operatorname{artanh}(1/2))=1/2.
\]
Every subsequent hidden layer applies
\[
 a'_1=\tanh(20a_1-20a_2-10c),\quad
 a'_2=\tanh(20a_1-20a_2-30c),\quad
 c'=\tanh(2\operatorname{artanh}(1/2)c)=1/2.
\]
If $2a_1-2a_2-2c=t$, the next two nonconstant neurons equal
$\tanh(10(t+1/2))$ and $\tanh(10(t-1/2))$.
The final readout $2a_1-2a_2-2c$ therefore gives $F^r(x_1)$.
Every stated weight is at most $30$ in absolute value, and the
constant rail uses only the input constant and internal linear weights.
\end{proof}

\section{A fixed input distribution and an approximation barrier}
\label{sec:distribution}
For independent fair signs $b_1,b_2,\ldots$, let
\[
 T=\lim_{j\to\infty}\phi_{b_1}\circ\cdots\circ\phi_{b_j}(0).
\]
The contraction bound in Lemma~\ref{lem:branches} makes this limit
well-defined. Let $\mu$ be its law and $K$ its compact support. Then
\[
 K\subset[-13/20,-1/2]\cup[1/2,13/20],\qquad F_*\mu=\mu.
\]
Define $x(t)=(t,\sqrt{1-t^2})$ on $[-13/20,13/20]$ and let
$D=x_*\mu$. This is a Borel probability distribution on $S^1$.
Since $\norm{x'(t)}_2=1/\sqrt{1-t^2}<2$ there, the embedding is
two-Lipschitz. In particular the same $D$ is used at every teacher depth.

\begin{lemma}[Small balls and sample separation]\label{lem:separation}
The measure $\mu$ is nonatomic. For every interval $I$ of length
$0<\ell\le1$,
\begin{equation}\label{eq:frostman}
 \mu(I)\le2\ell^\alpha,\qquad \alpha=\log2/\log20.
\end{equation}
For $n$ independent inputs from $D$, with probability at least $49/50$
their first coordinates are pairwise separated by at least $\Delta_n$
from~\eqref{eq:resources}.
\end{lemma}
\begin{proof}
Every sign word of length $j$ specifies a cylinder of mass $2^{-j}$.
Distinct depth-$j$ cylinders have distance at least $20^{-(j-1)}$:
their first differing branch supplies a gap of one, and their common
outer prefix contracts no distance by more than the lower bound in
\eqref{eq:branches}. The cylinder masses tending to zero show that
$\mu$ has no atoms.

Put $j=\lfloor\log_{20}(1/\ell)\rfloor$. If $j\ge1$, an interval
of length $\ell$ meets at most one depth-$j$ cylinder, because their
separation is strictly larger than $\ell$. Its mass is at most
$2^{-j}\le2\ell^\alpha$. The case $j=0$ follows from the trivial
mass bound one. Endpoints do not alter this argument.

Conditioning on one of two independent first coordinates gives
\[
 \Prb\{|T_i-T_j|<\Delta_n\}
 \le2(2\Delta_n)^\alpha.
\]
A union bound over unordered pairs is at most
$2n^2\Delta_n^\alpha=1/50$, since $2^\alpha\le2$.
\end{proof}

\begin{lemma}[Paired-input barrier, including posterior mixtures]
\label{lem:pairing}
Let $r\ge1$, $H\le3^r/10$, and $0\le h\le1/32$.
Suppose a distribution on predictors assigns mass at least $1-\beta$,
where $\beta\le1/3$, to functions within uniform distance $h$ on $S^1$
of some $H$-Lipschitz function. For target $f_r$ and test tolerance
$\tau=1/4$, its correct-output probabilities satisfy
\[
 D\{x:q(x)\le2/3\}\ge1/2.
\]
If the distribution is supported entirely on $H$-Lipschitz predictors,
the conclusion also holds with $1/2$ in place of the confidence
threshold $2/3$.
\end{lemma}
\begin{proof}
Let $w$ be a uniform sign word of length $r$, let $Y\sim\mu$ be
independent, and write $\phi_w$ for its composition. Pair
\[
 X_\pm=x(\phi_w(\phi_\pm(Y))).
\]
A fair choice between these two points has law $D$, while
\[
 |f_r(X_+)-f_r(X_-)|=|\phi_+(Y)-\phi_-(Y)|\ge1,
\]
\begin{equation}\label{eq:pair-gap}
 \norm{X_+-X_-}_2\le\frac{13}{5}\,3^{-r}.
\end{equation}
Thus an $H$-Lipschitz function changes by at most $13/50$ between
the pair. A predictor within $h$ of such a function cannot be correct
at both points: that would force the approximating function's change
to be at least $1-2(\tau+h)\ge7/16>13/50$.
For each pair, integration over the predictor distribution gives
\[
 q(X_+)+q(X_-)\le1+\beta\le4/3.
\]
At least one of the two probabilities is therefore at most $2/3$.
Averaging the balanced pairing proves the claim. If every predictor
is $H$-Lipschitz, the same sum is at most one and proves the final
assertion.
\end{proof}

\section{Exact memorization with a depth-independent Lipschitz ceiling}
\label{sec:memorization}
The next construction is useful because its energy grows almost linearly
with depth, while its product norm has a bound independent of depth.
It fits the observed real labels exactly.

\begin{lemma}[Exact bounded-label memorization]\label{lem:memorizer}
Let $d,n\ge1$, $L\ge3$, and suppose sample inputs have projected
coordinates $s_i=\langle u,x_i\rangle\in[-1,1]$, for some unit vector
$u$, with $|s_i-s_j|\ge\Delta>0$ whenever $i\ne j$.
For any labels $|y_i|\le7/8$, set
\[
 a=\Delta^{-1}\log(16n),\quad R=\max\{2,a\},\quad
 M=\left\lceil\max\{4,n,3dnR^2\}\right\rceil,\quad T=4Mn.
\]
At every width $N\ge M$ there is a depth-$L$ network with
\[
 z(x_i)=y_i,\qquad |z(x)|\le1\quad(x\in\R^d),\qquad
 \kappa_N(w)\le K:=2+(L+14)T^{1/(L-1)}.
\]
Furthermore
\begin{equation}\label{eq:product-bound}
 \left(\frac{K+1}{L}\right)^L\le e^{17}T^{3/2}.
\end{equation}
\end{lemma}
\begin{proof}
Order the coordinates as $s_1<\cdots<s_n$ and put
$c_j=(s_j+s_{j+1})/2$. For a vector $t\in\R^n$ define
\begin{equation}\label{eq:step-interpolator}
 g_t(x)=\frac{t_1+t_n}{2}+
 \sum_{j=1}^{n-1}\frac{t_{j+1}-t_j}{2}
       \tanh(a(\langle u,x\rangle-c_j)).
\end{equation}
The constant term is implemented using the hidden feature $1/2$ with
first-layer bias $\operatorname{artanh}(1/2)$ and readout weight
$t_1+t_n$. Thus the hidden width is $n$, including when $n=1$ and
the sum is empty.

Write $\theta_j=\tanh(a(\langle u,x\rangle-c_j))$. Since the
centers increase, $\theta_1\ge\cdots\ge\theta_{n-1}$.
Equation~\eqref{eq:step-interpolator} is a convex combination of the
$t_i$, with coefficients
\[
 \frac{1-\theta_1}{2},\quad
 \frac{\theta_1-\theta_2}{2},\quad\ldots,\quad
 \frac{\theta_{n-2}-\theta_{n-1}}2,\quad
 \frac{1+\theta_{n-1}}2.
\]
For $n=1$ the sole coefficient is one. Hence $|g_t|\le\norm t_\infty$
globally. Each sample is at distance at least $\Delta/2$ from every
threshold, so $|\tanh v-\operatorname{sign}v|\le2e^{-2|v|}$ gives
\[
 |g_t(x_i)-t_i|\le2(n-1)e^{-a\Delta}\norm t_\infty.
\]
Let $A$ be the linear map $t\mapsto(g_t(x_i))_{i=1}^n$. Then
\begin{equation}\label{eq:A-inverse}
 \norm{A-I}_{\infty\to\infty}<1/8,
 \qquad \norm{A^{-1}}_{\infty\to\infty}\le8/7.
\end{equation}

We next construct a scalar channel before selecting $t$.
Put $m=L-1\ge2$, $D_0=2T$ and
\[
 q=(D_0/4)^{1/(2m)}\ge1,\qquad
 k_s=\min\{s,m-s\},\qquad
 r_s=\frac{q^s}{\sqrt{D_0}\,k_s}\quad(1\le s\le m-1).
\]
These numbers satisfy $r_s\le1/(2qk_s)\le1/2$ and
$r_{m-1}=1/(2q)$. Define positive weights
\[
 \alpha_0=\operatorname{artanh}(r_1),\qquad
 b_s=\frac{\operatorname{artanh}(r_s)}{r_{s-1}}\quad(2\le s\le m-1),
\]
and the scalar function
\[
 G(t)=2q\tanh\bigl(b_{m-1}\tanh(\cdots\tanh(b_2\tanh(\alpha_0t))\cdots)\bigr).
\]
For $m=2$ this means $G(t)=2q\tanh(\alpha_0t)$.
By construction $G(1)=1$. It is odd and strictly increasing, and is
concave on $[0,1]$ because positive rescaling and composition preserve
increasing concave functions there. Thus $G$ maps $[-1,1]$ bijectively
onto itself, $G(t)\ge t$ on $[0,1]$, and
$|G^{-1}(y)|\le|y|$ for $|y|\le1$.

Choose
\[
 t=A^{-1}(G^{-1}(y_1),\ldots,G^{-1}(y_n))^T.
\]
Equation~\eqref{eq:A-inverse} gives
$\norm t_\infty\le(8/7)(7/8)=1$. Consequently $G\circ g_t$
has exact sample values $y_i$ and is globally bounded by one.

We verify its energy. For $\norm t_\infty\le1$, the first matrix
in~\eqref{eq:step-interpolator} has squared Frobenius norm at most
$3nR^2$. Its readout has squared norm at most $n+3\le4n$.
Pad with zero neurons to width $M$. The first-layer energy is then
at most one, and the weighted readout energy at most $T$.
For $N\ge M$, divide its $N$ neurons into $M$ balanced groups.
Duplicate each old neuron within its group, and divide its outgoing
readout coefficient by the group size. Since each group size is
between $N/(2M)$ and $2N/M$, this preserves the shallow function and
gives first energy at most two and weighted readout energy
$c^2=N\norm{v_N}_2^2\le2T=D_0$.

To implement $G$ after that first layer, make every row of the next
matrix equal to $\alpha_0v_N^T$. Each subsequent scalar-channel
matrix has all entries $b_s/N$, and the final readout has all entries
$2q/N$. Their respective energy costs are
$\alpha_0^2c^2\le D_0\alpha_0^2$, $b_s^2$, and $4q^2$.

The following elementary estimates bound their sum. The power series
for $\operatorname{artanh}$ gives
\begin{equation}\label{eq:atanh-bound}
 (\operatorname{artanh}(v)/v)^2\le1+v^2\quad(0\le v\le1/2).
\end{equation}
One verification is to bound
$\operatorname{artanh}(v)/v\le1+v^2/[3(1-v^2)]\le1+4v^2/9$;
squaring this upper bound gives at most $1+v^2$ for $v^2\le1/4$.
If $R_s=k_{s-1}/k_s$, then
\[
 \sum_{s=2}^{m-1}R_s^2\le m-2+7,\qquad
 R_s\le2,\qquad \sum_{s=1}^{m-1}k_s^{-2}<4.
\]
For the first inequality, pair the ratios $k/(k+1)$ and $(k+1)/k$
and use
\[
 \left(\frac{k+1}{k}\right)^2+
 \left(\frac{k}{k+1}\right)^2-2\le4/k^2,
\]
whose sum is less than seven. A possible central ratio equals one.
Equation~\eqref{eq:atanh-bound} now gives
\[
 D_0\alpha_0^2\le q^2+1/4,\qquad
 \sum_{s=2}^{m-1}b_s^2\le q^2(m-2+7)+4.
\]
The second remainder is bounded by
$\sum_s q^2R_s^2r_s^2\le\sum_s k_s^{-2}<4$.
Including a readout of size $q$ would cost at most $(m+12)q^2$;
the actual readout $2q$ adds $3q^2$. Including the first layer,
\[
 \kappa_N\le2+(m+15)q^2
 =2+(L+14)(T/2)^{1/(L-1)}\le K.
\]
Finally $K+1\le(L+17)T^{1/(L-1)}$ and $L/(L-1)\le3/2$,
which give~\eqref{eq:product-bound}.
\end{proof}

\section{Posterior energy and output ceilings}\label{sec:posterior}
\begin{lemma}[Strong-prior ceiling]\label{lem:strong-ceiling}
Fix a sample, a finite width $N$, depth $L\ge3$, and positive tolerance
$\rho$. Suppose it has an exact interpolant of energy at most $K$.
Then the posterior defined by~\eqref{eq:scales}--\eqref{eq:training}
satisfies
\[
 \Pi_{S,\sigma}\{\kappa_N>K+1\}\longrightarrow0
 \qquad(\sigma\downarrow0).
\]
If $K=2+(L+14)T^{1/(L-1)}$, it consequently assigns probability
tending to one to functions of Lipschitz constant at most
\begin{equation}\label{eq:Lip-ceiling}
 \Lambda=\frac{e^{17/2}}{\sqrt d}T^{3/4}.
\end{equation}
\end{lemma}
\begin{proof}
Continuity gives a parameter ball around the witness that still fits
the sample within $\rho$ and whose energies are at most $K+1/4$.
In the fixed finite parameter dimension its Gaussian probability has
lower exponential rate at least $-(K+1/4)/2$ at speed $\sigma^{-2}$.
The Gaussian radial tail outside energy $K+1$ has upper rate at most
$-(K+1)/2$. These bounds follow directly by integrating the density
over the ball and by the Gaussian squared-norm tail, respectively.
Their quotient tends to zero.

For an energy at most $K+1$, arithmetic--geometric mean and the
one-Lipschitz tanh activation give
\[
 \Lip(z_w)^2\le
 \norm{W^1}_F^2\cdots\norm{W^{L-1}}_F^2\norm v_2^2
 \le\frac1d\left(\frac{K+1}{L}\right)^L
 \le\frac{e^{17}}dT^{3/2},
\]
using~\eqref{eq:product-bound}. The constant input coordinate does
not change the Lipschitz bound with respect to the actual input.
\end{proof}

At finite exponents above one we need a ceiling for a nearby function,
rather than for the full parameter energy. We give its proof explicitly.

\begin{lemma}[Finite-scale approximation ceiling]\label{lem:finite-ceiling}
Suppose that for a sample on the Euclidean unit ball and every
$N\ge M$ there is an exact interpolant of energy at most
$K=2+(L+14)T^{1/(L-1)}$, where $L\ge3$ and $T\ge1$ are independent
of $N$. Fix $\zeta>0$, $0<\beta<1$, and $0<h<\rho$.
Define
\[
 A=\left(\frac{K+2}{L-1}\right)^{(L-1)/2},\qquad
 b^2=\max\{2(d+1)\log2+4\log8,\ 16\log9+8\log(8L)\},
\]
\[
 \eta=\min\left\{1,\frac1{4\sqrt{L(K+1)}},
                         \frac h{\sqrt2 LA}\right\}.
\]
For every
\begin{equation}\label{eq:width-threshold}
 N\ge\left\lceil\max\left\{M,2,(b/\eta)^{2/\zeta},
                         2\log(2/\beta)\right\}\right\rceil,
 \qquad 0<\sigma^2\le N^{-(1+\zeta)},
\end{equation}
the posterior assigns mass at least $1-\beta$ to outputs uniformly
within $h$ on the unit ball of a function with energy at most $K+1$
and Lipschitz constant at most~\eqref{eq:Lip-ceiling}.
The width threshold has a polynomial upper bound in
$d,M,T,L,h^{-1},\log(1/\beta)$ for fixed $\zeta$.
\end{lemma}
\begin{proof}
The energy-coordinate matrices are
\[
 A_1=\sqrt{d/N}W^1,\qquad A_j=W^j\ (2\le j<L),\qquad
 A_L=\sqrt N v^T.
\]
Their entries are independent $\mathcal N(0,\sigma^2)$ and their
total squared Euclidean norm is $\kappa_N$.
After normalizing the augmented input by $\sqrt d$ and hidden states
by $\sqrt N$, the first activation map is
$u\mapsto N^{-1/2}\tanh(\sqrt N A_1u)$, and the other hidden
activation maps have the same form. They are one-Lipschitz in their
linear argument and zero-preserving. The normalized input norm on
the unit ball is at most $\sqrt2$.

Suppose a residual network $U$ has energy at most $K+1$, and a
perturbation $B$ has operator norm at most $\eta$ in each normalized
block. Telescoping one layer at a time and applying
arithmetic--geometric mean to the other $L-1$ layer norms gives
\begin{equation}\label{eq:perturbation}
 \sup_{\norm x_2\le1}|z_{U+B}(x)-z_U(x)|
 \le\sqrt2 L\eta
 \left[\frac{K+1+2\eta\sqrt{L(K+1)}+L\eta^2}{L-1}
 \right]^{(L-1)/2}.
\end{equation}
Indeed, writing the residual block norms as $a_j$, each telescoping
term is at most $\sqrt2\eta$ times a product of $L-1$ numbers
$a_j+\eta$. Their squared sum is bounded by the numerator, since
$\sum a_j^2\le K+1$ and $\sum a_j\le\sqrt{L(K+1)}$.
The choice of $\eta$ makes the cross term at most $1/2$ and
$L\eta^2\le1/[16(K+1)]$. Thus that numerator is at most $K+2$
and~\eqref{eq:perturbation} is at most $\sqrt2L\eta A\le h$.

Put $t^2=N\sigma^2$. Let
\[
 a_d^2=2(d+1)\log2+4\log8,\qquad
 b_L^2=16\log9+8\log(8L).
\]
In energy coordinates let $\mathcal T$ be the set where first and
last block Frobenius norms are at most $ta_d$, and internal block
operator norms are at most $tb_L$. This set is nonempty, closed,
convex and centrally symmetric, and its Gaussian probability is
at least $1/2$. To verify the last assertion, the first squared norm
has law $(t^2/N)\chi^2_{N(d+1)}$; the last uses $N$ degrees of
freedom. The exponential moment at parameter $1/4$ bounds either
failure probability by $1/8$. A $1/4$-net of the unit sphere has at
most $9^N$ points. Applying the scalar Gaussian tail on two such
nets gives, for each internal matrix $tG/\sqrt N$,
\[
 \Prb\{\norm{tG/\sqrt N}_{\rm op}>tb_L\}
 \le2\exp(2N\log9-Nb_L^2/8)\le1/(4L).
\]
The union bound proves the required mass. Every normalized block
in $\mathcal T$ has operator norm at most $tb$.

Condition~\eqref{eq:width-threshold} gives $t\le N^{-\zeta/2}$,
so $tb\le\eta$. If $w_S$ is the exact-fitting witness, then
$w_S+\mathcal T$ fits the sample within $h<\rho$ by the perturbation
bound. The Gaussian shift identity and central symmetry give
\[
 \mu(w_S+\mathcal T)
 =e^{-\kappa_N(w_S)/(2\sigma^2)}
   \int_{\mathcal T}e^{-\langle w_S,B\rangle_\kappa/\sigma^2}\,d\mu(B)
 \ge\tfrac12e^{-K/(2\sigma^2)}.
\]
Pairing $B$ with $-B$ in the integral proves the inequality.

Gaussian isoperimetry~\cite{borell,ledoux}, applied to a set of
Gaussian mass at least $1/2$, gives
\[
 \mu\{\operatorname{dist}_\kappa(w,\mathcal T)^2>K+1\}
 \le e^{-(K+1)/(2\sigma^2)}.
\]
Dividing by the preceding lower bound for the conditioning event
shows that the posterior probability of this event is at most
$2e^{-1/(2\sigma^2)}\le\beta$. Off this event, let $B$ be the metric
projection of $w$ onto $\mathcal T$ and set $U=w-B$.
Then $\kappa_N(U)\le K+1$, and~\eqref{eq:perturbation} gives
$\sup|z_w-z_U|\le h$. The product-norm bound gives the stated
Lipschitz ceiling. No claim that $w$ itself has low energy is needed.

Finally
\[
 K+2\le(L+18)T^{1/(L-1)},\qquad
 A^2\le\left(\frac{L+18}{L-1}\right)^{L-1}T\le e^{19}T.
\]
Thus every inverse radius and threshold in the lemma has the claimed
polynomial upper bound for fixed $\zeta$.
\end{proof}


\section{Proof of the exponential lower bounds}\label{sec:proof}
For fixed $d=2$, the quantities in~\eqref{eq:resources}--\eqref{eq:MT}
obey
\[
 T_n\le C_0 n^{4+4/\alpha},\qquad
 \Lambda_n:=\frac{e^{17/2}}{\sqrt2}T_n^{3/4}
 \le C_1 n^{3+3/\alpha}
\]
for universal constants $C_0,C_1$. For instance use
$\log(16n)\le Cn$ in the displayed definitions. The exponent
$3+3/\alpha$ is less than $16$. Hence there is a universal $r_0\ge2$
such that, for all $r\ge r_0$ and $n\le e^{r/100}$,
\begin{equation}\label{eq:exp-comparison}
 \Lambda_n\le3^r/10.
\end{equation}
No constant in this comparison depends on learner width or depth.

\begin{proof}[Proof of Theorem~\ref{thm:strong}]
The distribution and teachers were constructed in
Sections~\ref{sec:map}--\ref{sec:distribution}.
On the sample-separation event of Lemma~\ref{lem:separation}, all labels
belong to $K\subset[-13/20,13/20]$, so
Lemma~\ref{lem:memorizer} supplies an exact-fitting witness with energy
at most $K_n(L)$ at every stated width. For each fixed such sample and
fixed width, Lemma~\ref{lem:strong-ceiling} makes the posterior mass
outside the $\Lambda_n$-Lipschitz class tend to zero. By
\eqref{eq:exp-comparison}, Lemma~\ref{lem:pairing} applies as soon as
that exceptional mass is at most $1/3$. Thus the indicator of the
sample event inside~\eqref{eq:strong-result} is eventually one for
each separated sample. Fatou's lemma and the probability $49/50$ of
separation prove the asserted lower limit. This use of Fatou does not
require a small-noise threshold uniform over realized samples.
\end{proof}

\begin{proof}[Proof of Theorem~\ref{thm:finite}]
Use Lemma~\ref{lem:finite-ceiling} with $K=K_n(L)$, $M=M_n$,
$\beta=1/4$ and $h=\rho/4=1/32$, increasing its width threshold to
at least $30$. On every separated sample the exact memorizer exists
and the lemma gives posterior mass at least $3/4$ on outputs within
$h$ of $\Lambda_n$-Lipschitz functions. Equations
\eqref{eq:exp-comparison} and Lemma~\ref{lem:pairing} give the
confidence-error conclusion. Separation has probability at least
$49/50$. The explicit width threshold has a polynomial upper bound
in $n,L$ for fixed $\zeta$, as proved in that lemma.
\end{proof}

\section{Exact interpolation and the order of limits}\label{sec:exact}
\begin{corollary}[Exact minimum-energy interpolants also fail]
\label{cor:exact}
Under the parameters of Theorem~\ref{thm:strong}, with sample
probability at least $49/50$, every minimum-energy exact interpolant
at every fixed width $N\ge\max\{30,M_n\}$ is $\Lambda_n$-Lipschitz
and has regression error at least $1/2$ at tolerance $1/4$.
Every mixture of such interpolants has confidence-error mass at least
$1/2$, at confidence threshold $2/3$ or even $1/2$.

At each such finite width, iterated weak limits of the posterior
obtained by first sending $\sigma\downarrow0$ at fixed positive
training tolerance and then sending that tolerance to zero exist
subsequentially. Every such limit is supported on minimum-energy exact
interpolants and obeys the same lower bounds.
\end{corollary}
\begin{proof}
For a separated sample the exact memorizer has energy at most $K_n(L)$.
The exact feasible set is nonempty and closed, and the quadratic energy
is coercive. Its minimum therefore exists. Every minimizer has energy
at most $K_n(L)$, so the arithmetic--geometric mean calculation in
Lemma~\ref{lem:strong-ceiling} bounds its Lipschitz constant by
$\Lambda_n$. The pairing argument gives error at least $1/2$ for
each function and the stated conclusions for all mixtures.

For completeness, fix the sample and width, and put
\[
 e_S(w)=\max_i|z_w(X_i)-Y_i|,\quad
 b_a=\inf_{\{e_S<a\}}\kappa_N,\quad
 b_0=\min_{\{e_S=0\}}\kappa_N.
\]
An open ball within $\{e_S<a\}$ of energy arbitrarily close to
$b_a$ and the Gaussian tail estimate show that, as $\sigma\downarrow0$,
posterior energy concentrates at $b_a$. The measures are tight and
every weak limit is supported on
$\{e_S\le a,\ \kappa_N=b_a\}$. In particular it is supported on
the common compact set $\{\kappa_N\le b_0\}$.

As $a\downarrow0$, the numbers $b_a$ increase to $b_0$. To verify
the nontrivial inequality, take approximate minimizers at tolerances
tending to zero. They lie in a compact energy sublevel, and a
convergent subsequence has exact fit with energy at most
$\lim_{a\downarrow0}b_a$. Every outer weak limit is consequently
supported on $\{e_S=0,\kappa_N=b_0\}$. Compactness supplies the
subsequential limits at both stages.
\end{proof}

The order of limits in the corollary is part of the statement. It does
not identify a regular conditional posterior given exact real-valued
observations at a fixed prior scale, or describe arbitrary simultaneous
limits of the scale and tolerance.

\section{An exponential lower bound at the mean-field exponent}
\label{sec:critical-depth-lower}

At the critical exponent $\gamma=1$, a different argument gives an
exponential lower bound for the same teacher family and the same fixed
nonatomic input distribution $D$. The width tends to infinity with
teacher depth and the realized sample fixed. Only finite-dimensional
large deviations are needed: apply them separately to the training
sample and to one test pair, then integrate over the test pairing.

\begin{theorem}[Critical regression lower bound]
\label{thm:critical-regression-depth-lower}
Use the distribution $D$ and targets $f_h(x)=F^h(x_1)$ from
Theorem~\ref{thm:strong}. There are universal constants $c>0$ and $h_0$
such that the following holds for every $h\ge h_0$ and integer
$1\le n\le e^{ch}$. Draw $n$ independent inputs from $D$ and observe
their exact teacher outputs. Use the original Gaussian prior at depth
$L=h+1$, with first-layer entry variance $1/2$, internal entry variance
$1/N$, and readout entry variance $1/N^2$. Condition solely on
\[
 |z_w(X_i)-f_h(X_i)|<\rho,\qquad \rho=1/8.
\]
With sample probability at least $49/50$, the posterior satisfies
\[
 \liminf_{N\to\infty}
 \mathbb E_{w\mid S}D\{|z_w-f_h|>1/4\}\ge\frac12.
\]
Moreover, writing
$p_N(x)=\Pr_{w\mid S}\{|z_w(x)-f_h(x)|\le1/4\}$, the same samples obey
\[
 \liminf_{N\to\infty}D\{x:p_N(x)<2/3\}\ge\frac12.
\]
The teacher still has width three and weights bounded by thirty.
There is no global output-margin condition. These statements concern
the actual posterior, including mixtures, and give no width threshold
uniform in growing $h$.
\end{theorem}

\subsection{Finite-domain rates and a metric bound}

The rate below measures the exponential prior cost of approximating
prescribed outputs on finitely many inputs. Its hidden-layer terms are
relative entropies, and its final term is the Gaussian readout cost.
We first show that a small total cost limits the output difference
between nearby inputs.

We include the large-deviation argument needed below. It works for any
fixed finite input set $X=\{x^1,\ldots,x^m\}\subset S^1$. For a feature
law $\nu$ on $E=[-1,1]^m$, write
$Q(\nu)=\int hh^T\,d\nu(h)$. Let $\mathsf M_Q$ be the law of
coordinatewise tanh applied to $\mathcal N(0,Q)$, and put
\[
 Q_0(x,y)=\frac{1+x\cdot y}{2},\qquad
 \mathcal G_Q(z)=
 \begin{cases}\frac12z^TQ^+z,&z\in\operatorname{range}Q,\\
 +\infty,&\text{otherwise}.
 \end{cases}
\]
Here and below singular covariance matrices are allowed.

\begin{lemma}[Finite-domain critical large deviations]
\label{lem:critical-finite-ldp-depth}
The output vector $Z_N=(z_w(x))_{x\in X}$ obeys a large-deviation
principle with speed $N$ and good rate
\[
 I_L(z)=\inf_{\nu_1,\ldots,\nu_{L-1}}
 \left\{\sum_{s=1}^{L-1}D(\nu_s\Vert\mathsf M_{Q_{s-1}})
                   +\mathcal G_{Q_{L-1}}(z)\right\},
 \qquad Q_s=Q(\nu_s).
\]
For every finite-cost path of total cost $A$ in this formula,
\begin{equation}\label{eq:critical-metric-bound-depth}
 |z(x)-z(y)|\le\sqrt A\,|x-y|
                         \exp\!\bigl(\sqrt{2(L-1)A}\bigr)
 \qquad(x,y\in X).
\end{equation}
\end{lemma}

\begin{proof}
Let $\Gamma_m=\mathcal N(0,I_m)$, and take independent empirical
measures $\eta_{s,N}$ of $N$ samples from $\Gamma_m$, for
$1\le s<L$. Independently take $V_N=G/\sqrt N$, $G\sim\Gamma_m$.
For $\mathsf T_Q(u)=\tanh(Q^{1/2}u)$, recursively define
\[
 \nu_{1,N}=(\mathsf T_{Q_0})_\#\eta_{1,N},\qquad
 \nu_{s,N}=(\mathsf T_{Q(\nu_{s-1,N})})_\#\eta_{s,N}.
\]
These are exactly the network's empirical hidden-feature laws on $X$:
conditionally on the preceding layer, its next rows are independent
Gaussian fields with that empirical second moment. The final output
has exactly the law $Q(\nu_{L-1,N})^{1/2}V_N$.

Sanov's theorem and the Gaussian large-deviation principle give joint
good rate $\sum_s D(\eta_s\Vert\Gamma_m)+\|v\|^2/2$.
The recursive map just described is continuous in the weak topology
for the measures and the ordinary topology for $v$. Indeed, matrix
square roots are continuous; the bounded continuous test functions
composed with $\mathsf T_Q$ converge uniformly on compact sets when
$Q$ varies; and weakly convergent measures form a tight sequence.
Second moments on the compact feature space $E$ are also continuous.
The contraction principle therefore applies. The entropy-contraction
identity
\[
 \inf_{\eta:T_\#\eta=\nu}D(\eta\Vert\Gamma_m)
       =D(\nu\Vert T_\#\Gamma_m)
\]
follows by data processing and, for the reverse inequality, the lift
$d\eta=q(T(u))\,d\Gamma_m(u)$ when
$q=d\nu/d(T_\#\Gamma_m)$. Finally
$\inf_{Q^{1/2}v=z}\|v\|^2/2=\mathcal G_Q(z)$.
These identities give the displayed rate, including the singular cases.
This is an application of the standard Sanov and contraction theorems
\cite{dembozeitouni}.

For the metric bound, write $D_s$ for each hidden-layer entropy and
$r$ for the readout cost. Inverse tanh lifts a finite-entropy feature
law to a law $\eta$ with the same entropy relative to its reference
Gaussian. The Gaussian quadratic transport inequality
\cite{talagrand1996}, applied after whitening on the covariance range,
gives a coupling $(U,G)$ satisfying
\[
 \mathbb E\|Q^{+1/2}(U-G)\|^2\le2D_s.
\]
For $b=e_x-e_y$, Cauchy--Schwarz and Minkowski imply
\[
 \|b^TU\|_{L^2}\le(1+\sqrt{2D_s})\sqrt{b^TQb}.
\]
Tanh is $1$-Lipschitz, so this is also an upper bound on the next
feature distance. The initial distance is $|x-y|/\sqrt2$.
The readout inequality is
$|b^Tz|\le\sqrt{2r}\sqrt{b^TQ_{L-1}b}$.
Multiply these inequalities and use
\[
 \sum_{s=1}^{L-1}\log(1+\sqrt{2D_s})
 \le\sqrt{2(L-1)\sum_sD_s}\le\sqrt{2(L-1)A},\qquad r\le A.
\]
This proves~\eqref{eq:critical-metric-bound-depth}.
\end{proof}

Each test pair $(X_+,X_-)$ from Lemma~\ref{lem:pairing} has
input distance at most $3\cdot3^{-h}$ and target gap at least one.
Predicting both to tolerance $1/4$ requires an output gap of at least
$1/2$. At $L=h+1$, the metric bound shows that every such output on
the two-point domain has
\begin{equation}\label{eq:critical-correct-pair-rate}
 I_L(z)\ge h/100
\end{equation}
for all sufficiently large $h$: a path of cost $A\le h/100$ would
instead give gap at most
$3\cdot3^{-h}\sqrt{h/100}\exp(h/\sqrt{50})$, which tends to zero.
The same conclusion holds after taking the infimum over paths.

\subsection{A low-rate training interpolator}

The training constraints, in contrast, can be met at a much smaller
rate. The next construction spends entropy only in the first layer and
in logarithmically many final layers. Most hidden layers follow the
uncontrolled Gaussian recursion, which reduces their variances while
retaining a weak signal correlated with the labels. The final controlled
layers amplify that signal.

\begin{lemma}[A first-feature mean seed]
\label{lem:critical-mean-seed-depth}
Suppose $t_1<\cdots<t_n$ belong to $[-2/3,2/3]$ and have
separation at least $\Delta\le1$. For any signs $y_i$ and any fixed
finite domain containing $x(t_i)$, a genuine first-feature law $\nu$
satisfies
\[
 D(\nu\Vert\mathsf M_{Q_0})\le D,\qquad
 \mathbb E_\nu H(x(t_i))=m_n y_i,\qquad m_n=\frac1{2n},
\]
where $1\le D\le C\Delta^{-2}[1+\log(n+1)]^2$ for a universal $C$.
\end{lemma}

\begin{proof}
Use the ideal step vectors $s_0=1$ and
$s_j(i)=\operatorname{sign}(t_i-(t_j+t_{j+1})/2)$, $1\le j<n$.
The inverse of this basis maps the maximum norm to coefficient
$\ell^1$ norm with norm at most $n$: the coefficients are successive
half-differences and the half-sum of the two endpoint values.

In first-layer weight space, shift the bias and first coordinate so
that a row's preactivation mean is $B(t-u_j)$, where $u_j$ is the
corresponding midpoint. For the constant vector use constant mean $B$.
The random preactivation variance is one on $S^1$. Since
\[
 1-\mathbb E\tanh(m+Z)\le2e^{-2m+2}\qquad(Z\sim\mathcal N(0,1)),
\]
choosing $B=C\Delta^{-1}[1+\log(n+1)]$ makes the resulting mean
vectors $v_j$ approximate $s_j$ to maximum error $1/(4n)$.
A Neumann series gives exact coefficients
$y=\sum_j c_jv_j$ with $\sum_j|c_j|\le4n/3$.
Mix the corresponding signed Gaussian shifts with weights
$m_n|c_j|$, and put the remaining probability on the unshifted
Gaussian law. Negating a shift negates its feature mean by oddness.
The desired mean identity follows. Each shift has Gaussian relative
entropy at most $CB^2$, and convexity and data processing give the
stated bound. These are laws of actual weight vectors, so their
feature means and covariances are defined on the entire finite domain,
not only on the sample.
\end{proof}

\begin{lemma}[Free cooling and late amplification]
\label{lem:critical-late-amplification-depth}
Suppose a first-feature law has entropy at most $D\ge1$ and mean
$a_0y_i$ at selected inputs, where $y_i\in\{-1,1\}$ and
$0<a_0\le1$. For every fixed $0<\eta<1$, there are constants
$C_1,C_2$, depending only on $\eta$, such that if
\[
 L\ge C_1[1+\log(D/a_0^2)+\log(L+1)],
\]
then the critical rate admits an output $z$ with
\[
 |z(x_i)-a y_i|<a\eta,
 \qquad I_L(z)\le C_2[1+\log(D/a_0^2)+\log(L+1)],
 \qquad a=3/5.
\]
\end{lemma}

\begin{proof}
Mix the supplied law with the reference first-feature law at weight
$1/D$. This costs at most one and yields a covariance $Q\succeq bb^T$
with $b(x_i)=a_0y_i/\sqrt D$. The vector $b$ is the full-domain mean
of the supplied law divided by $\sqrt D$.

Consider an uncontrolled transition with $Q\succeq bb^T$ and
$q_x=Q(x,x)\le1$. Write the reference Gaussian field as $bZ_0+G_R$.
Project its tanh features onto $Z_0$. Gaussian integration by parts
shows that the resulting covariance dominates $b'b'^T$, where
\[
 b'_x=b_xh(q_x),\qquad
 h(q)=\mathbb E\operatorname{sech}^2(\sqrt qZ)=1-F_0(q),\qquad
 F_0(q)=\mathbb E\tanh^2(\sqrt qZ).
\]
The new diagonal is $F_0(q_x)$. There are universal $c>0,C<\infty$
such that
\[
 F_0(q)\le q-cq^2,\qquad -\log h(q)\le Cq\quad(0\le q\le1).
\]
For example $c=1/100$ works. On $|u|\le1$,
$u^2-\tanh^2u\ge u^4/6$; integrate over $|Z|\le1$ to obtain the
first bound. For the second, $h(q)\ge h(1)\ge1/8$ and $F_0(q)\le q$.
It follows that after $H$ free layers the diagonal is at most
$(1+cH)^{-1}$ and the signed seed at every selected input is at least
\begin{equation}\label{eq:critical-free-seed}
 u_*:=e^{-8}(L+1)^{-800}a_0/\sqrt D
 \qquad(H\le L).
\end{equation}
Indeed the diagonal reciprocals grow by at least $c$ per step, and
summing $-\log h(q)\le8q$ proves this loose explicit bound.

Now suppose $Q\succeq bb^T$ and $y_ib_i\ge u>0$ on the sample.
Shift the reference Gaussian by $A b$. The entropy cost is at most
$A^2/2$, since $b^TQ^+b\le1$. The new mean $b'$ satisfies
\[
 y_i b'_i\ge c_0\min\{Au,1\},\qquad c_0=1/100,
\]
and the new covariance dominates $b'b'^T$. To check the inequality,
integrate the derivative
$\mathbb E\operatorname{sech}^2(v+\sqrt qZ)$ for $0\le v\le1$
and $q\le1$, restricting to $|Z|\le1$; larger means follow by
monotonicity. Choose $A=2/c_0$. Repeating this shift doubles the
minimum signed mean until it reaches $c_0/2$, and preserves this
lower bound thereafter.

Choose
\[
 m=\max\{0,\lceil\log_2(c_0/(2u_*))\rceil\},\qquad
 H=L-m-3.
\]
The stated depth condition, with a sufficiently large $C_1$, ensures
$H\ge0$. Use one initial seed layer, $H$ free layers, $m$ shift
layers, and one final saturation layer. This is exactly $L-1$ hidden
layers and avoids any circular choice of the cooling time.
For the final shift use $B b$, with $B$ large enough in terms of
$\eta$ that its feature mean $v$ obeys $|v_i-y_i|<\eta$.
The bound
\[
 1-y_i\mathbb E\tanh(G_i+B b_i)\le2e^{2-Bc_0}
\]
suffices. Its entropy cost is at most $B^2/2$. Finally $Q_{L-1}\succeq vv^T$,
so $z=av$ has readout cost at most $a^2/2$.
The total cost is at most $1+mA^2/2+B^2/2+a^2/2$, which has the
asserted bound. Every Gaussian shift and rank-one domination is on
the full finite domain throughout the construction.
\end{proof}

\subsection{Comparison of posterior masses}

Correct prediction at both members of a test pair has a larger prior
rate than fitting the training sample. Dividing the former probability
by the latter therefore makes simultaneous correctness on a fixed pair
vanish as width grows. Averaging over the pairing then gives the risk
and confidence bounds for the full posterior.

\begin{proof}[Proof of Theorem~\ref{thm:critical-regression-depth-lower}]
By Lemma~\ref{lem:separation}, the training scalar coordinates are
$\Delta_n$-separated with probability at least $49/50$, where
$\Delta_n=(100n^2)^{-1/\alpha}$ and $\alpha=\log2/\log20$.
Fix any such sample. Since the Cantor measure is invariant under $F$,
every teacher label belongs to
$[-13/20,-1/2]\cup[1/2,13/20]$.
Set $y_i=\operatorname{sign}f_h(X_i)$ and $a=3/5$. Thus
\[
 |f_h(X_i)-ay_i|\le1/10.
\]
Apply Lemma~\ref{lem:critical-mean-seed-depth} on the finite training
set. Its quantity $D/m_n^2$ is polynomial in $n$.
Lemma~\ref{lem:critical-late-amplification-depth}, with $\eta=1/100$,
then supplies an output whose training errors are less than
\[
 1/10+a\eta=53/500<1/8.
\]
Reduce the universal $c$ and increase $h_0$ so that, whenever
$n\le e^{ch}$, the depth condition holds and this witness's rate is
less than $h/200$. The open training event $A_S$ therefore satisfies
\begin{equation}\label{eq:critical-training-rate-depth}
 \liminf_{N\to\infty}\frac1N\log\Pr(A_S)>-h/200.
\end{equation}
The strict inequality follows from the strict witness cost bound.

Now fix one test pair $(X_+,X_-)$ in the balanced pairing of
Lemma~\ref{lem:pairing}. Let $B_{X_+,X_-}$ be the closed event that
both real outputs are within $1/4$ of their teacher values.
Apply the two-point version of Lemma~\ref{lem:critical-finite-ldp-depth}.
Equation~\eqref{eq:critical-correct-pair-rate} gives
\[
 \limsup_{N\to\infty}\frac1N\log\Pr(B_{X_+,X_-})\le-h/100.
\]
Together with~\eqref{eq:critical-training-rate-depth}, the elementary
ratio bound
\[
 \Pr_{w\mid S}(B_{X_+,X_-})
 \le\frac{\Pr(B_{X_+,X_-})}{\Pr(A_S)}
\]
shows that this posterior probability tends to zero for every fixed
test pair. No joint or functional large-deviation principle is being
assumed. The width at which the ratio becomes small may depend on
the test pair.

Let $b_N(X_+,X_-)$ denote this pair-correct probability and average
with respect to the pairing law. Bounded convergence gives
$\mathbb E b_N\to0$. Inclusion--exclusion gives, pair by pair,
\[
 p_N(X_+)+p_N(X_-)\le1+b_N(X_+,X_-).
\]
A fair choice between the paired inputs has law $D$. Averaging proves
that the posterior Gibbs error is at least
$(1-\mathbb E b_N)/2$, giving its asserted lower limit.

For the confidence assertion, if $b_N<1/3$, both correct-output
probabilities cannot be at least $2/3$. At least one member of the
pair has probability strictly below $2/3$. Consequently
\[
 D\{x:p_N(x)<2/3\}\ge\frac12\Pr_{\rm pair}\{b_N<1/3\}
 \ge\frac{1-3\mathbb E b_N}{2}.
\]
Taking lower limits proves the confidence claim, since $\mathbb E b_N\to0$.
The argument applies
for every separated sample and handles arbitrary posterior mixtures.
\end{proof}


\section{Smooth input densities when the distribution may vary with depth}
\label{sec:smoothing}
\begin{corollary}[Smooth-distribution variant]\label{cor:smooth}
For every $r$ there is a probability distribution $D_r$ on $S^1$
with a smooth density with respect to arclength such that
Theorems~\ref{thm:strong} and~\ref{thm:finite} and
Corollary~\ref{cor:exact} hold with $D_r$ replacing $D$.
The teacher $f_r$, sample range, resource bounds and confidence-error
conclusions are unchanged. The width-first mean-field conclusion of
Theorem~\ref{thm:critical-regression-depth-lower} also holds, with
universal constants of the same kind, which may be chosen to serve both
statements.
No uniform bound on the density or its
derivatives is asserted as $r$ grows.
\end{corollary}
\begin{proof}
Let $U_r$ have a smooth compactly supported probability density on
$[-10^{-2}20^{-r},10^{-2}20^{-r}]$, independently of $T\sim\mu$,
and let $D_r$ be the law of $x(T+U_r)$. The scalar law is the
convolution of $\mu$ with that smooth density. Its support is in
$(-.66,.66)$, so its pushforward has a smooth arclength density on
the circle, extended by zero off this arc.

For every interval $I$, conditioning on $U_r$ preserves the bound
$\Prb(T+U_r\in I)\le2|I|^\alpha$. Thus the sample-separation
probability and all resources remain unchanged. Globally $|F'|\le20$,
so
\[
 |F^r(t+u)-F^r(t)|\le20^r|u|\le1/100.
\]
The real training labels therefore have magnitude at most
$13/20+1/100=.66<7/8$, as required by
Lemma~\ref{lem:memorizer}.

Use the same jitter $U_r$ for both members of each pair in
Lemma~\ref{lem:pairing}. The scalar input gap is unchanged, the arc
embedding remains two-Lipschitz, and the target gap is at least
$1-2/100=.98$. For the test tolerance and every $h\le1/32$, the
lower bound on an approximating function's change is now
$.98-2(1/4+h)\ge.4175$, still greater than $13/50$.
The same pairing and posterior arguments therefore apply.

For the mean-field statement, the smoothed real label differs from
$3/5$ times its sign by at most $.11$. The low-rate witness in
Theorem~\ref{thm:critical-regression-depth-lower} adds less than
$.006$ training error, so its total error is less than $.116<1/8$.
The first-feature seed applies on the slightly enlarged arc and
uses the unchanged sample separation. Correct prediction of both
members of a smoothed pair requires an output gap of at least
$.98-2(1/4)=.48$. The same entropy-based metric bound excludes this
at rate at most $r/100$ for large $r$. The finite-pair posterior
ratio and bounded-convergence argument then prove the mean-field
conclusion as well.
\end{proof}

\section{Scope and interpretation}
The construction rules out a general guarantee jointly polynomial in
teacher width, bounded weights and depth for the stated Gaussian learners
on arbitrary continuous-domain distributions. Given any proposed joint
polynomial, one can choose depth large enough that its sample requirement
is below $e^{(L-1)/100}$ and then choose an allowed learner width from
the theorem. In the strong-prior case the failure persists arbitrarily
far into the small-scale limit. At every fixed exponent above one it
already occurs at the actual finite prior scale, for sufficiently large
polynomial learner width.

The distributional qualification is substantial. The fixed input law is
nonatomic but singular: it is supported on a Cantor subset of a circle
arc. Smooth arclength densities are possible when the input law may
vary with depth, by Corollary~\ref{cor:smooth}. This paper does not
prove the same lower bound for uniform measure on the circle or sphere,
for one fixed smooth input density, or for every smooth density. Nor does it
claim a lower bound for every possible learning algorithm. Its proof
uses the specific energy competition induced by the Gaussian prior.

The continuous domain matters. It contains distinct points at
exponentially small Euclidean distances. Distinct Boolean-cube vertices
have distance at least two, so this close-pair mechanism does not transfer
to the original Boolean classification model. A coordinate restriction
or a grid encoding by Boolean strings also changes the network
representation and its norm. Consequently this result leaves open the
general small-deep-circuit question on the Boolean cube.

The Gaussian-distance proof requires $\gamma>1$ because the typical
Gaussian component then has operator norm tending to zero with width.
The separate mean-field proof instead uses finite-dimensional entropy
costs and takes width to infinity without a quantified threshold.
At this endpoint, the positive sample exponent need not equal
$1/100$, and the one-half confidence-error bound is a lower limit
in width rather than a finite-width assertion.
The theorem concerns increasing \emph{teacher} depth, and does not
contradict polynomial dependence on additional learner depth for a fixed
shallow target. No claim about non-Gaussian priors is made here.

\begingroup\small
\begin{thebibliography}{9}
\bibitem{borell} C. Borell. The Brunn--Minkowski inequality in Gauss space.
\emph{Inventiones Mathematicae} 30 (1975), 207--216.
\bibitem{ledoux} M. Ledoux. \emph{The Concentration of Measure Phenomenon}.
American Mathematical Society, 2001.
\bibitem{dembozeitouni} A. Dembo and O. Zeitouni.
\emph{Large Deviations Techniques and Applications}. Second edition,
Springer, 1998.
\bibitem{talagrand1996} M. Talagrand. Transportation cost for Gaussian
and other product measures. \emph{Geometric and Functional Analysis}
6 (1996), 587--600.
\end{thebibliography}\endgroup
\end{document}
